% Source: MQ18b_v1.html

\documentclass[11pt]{article}
\usepackage[margin=1in]{geometry}
\usepackage{amsmath}
\usepackage{amssymb}
\usepackage{siunitx}
\usepackage{enumitem}

\title{MQ18b: Heat and Calorimetry}
\author{}
\date{}

\begin{document}
\maketitle

\section*{MQ18b: Heat and Calorimetry}

\begin{enumerate}[leftmargin=*,label=\textbf{Question \arabic*.}]

\item \hfill \textit{1 pts}

The average translational kinetic energy of a single atom is 5.5 × 10$^{-22}$J in a container of helium. Calculate the temperature of the gas as a whole.

\item \hfill \textit{1 pts}

Suppose 600,796 atoms of an ideal gas are introduced into an evacuated bottle of volume 47 cm$^{3}$. Calculate the translational kinetic energy of the gas at a temperature of 552 K. Answer in fmJ (femtojoules). (10$^{15}$ fmJ= 1 J)



\emph{K}$_{tr}$ = \rule{2.5cm}{0.4pt} fJ

\item \hfill \textit{1 pts}

Calculate the most probable speed (in m/s) of molecular oxygen at a temperature of 97.7 °C. (Molar mass of atomic oxygen is 16.0 g/mol.)

\item \hfill \textit{1 pts}

Suppose the root-mean-square speed of a gas is 628 m/s. If the molecular weight in grams per mole of the gas is 10 g/mol, calculate the temperature of the gas in Kelvin.

\item \hfill \textit{1 pts}

A gas has 4 degrees of freedom. By the equipartition of energy, calculate the internal energy change in kJ required to raise the temperature of 1.5 moles of the gas by 54 C°.

\item \hfill \textit{1 pts}

Calculate the total internal energy in kJ of 6.6 mol of helium gas at a temperature of 25°C.



\emph{E}= \rule{2.5cm}{0.4pt} kJ

\end{enumerate}

\end{document}
